% Draft (Agent F) of the new front part of paper.tex, version 2. Merged into the top-level paper.tex at integration.

\section{The question}

I.~J.~Halvorsen and S.~Skogestad, \emph{Minimum energy consumption in multicomponent distillation. 3. More than three
products and generalized Petlyuk arrangements}, Ind.\ Eng.\ Chem.\ Res.\ 42 (2003) 616--629~\autocite{HS3}. The
abstract states:
\begin{openquote}
``We also conjecture that the minimum energy requirement for the generalized extended Petlyuk arrangement is lower than
the minimum-energy requirement for any distillation configuration when we consider conventional adiabatic sections and
no internal heat exchange.''
\end{openquote}
Section~6.3 (p.~626) makes the assumptions explicit and records the status:
\begin{openquote}
``We assume constant relative volatilities, constant molar flows, constant pressure, and no internal heat integration.
Then the generalized adiabatic Petlyuk arrangement has the lowest need for vaporization compared to any other adiabatic
distillation arrangement for separation of an arbitrary feed mixture into its pure components. We have not presented a
complete proof, so the above conclusion is a conjecture.''
\end{openquote}
The number of stages is infinite; an adiabatic section is ``a column section with constant molar flows and no heat
exchange along the section'' (\autocite{HS3}, \S6). Even the minimum vapor of the Petlyuk arrangement itself was left
partly open (Section~5.3, p.~624):
\begin{openquote}
``We have not carried out a detailed proof for the general N-component M-product case for other than the far left and
right peaks.''
\end{openquote}
The conjecture is restated in~\autocite{HS02}, where infinitely many stages are assumed ``because this provides a lower
bound on the energy requirement for ideal mixtures'', the energy is measured by the total vaporization, and the case
with internal heat exchange is called unknown. For three components, Fidkowski and Kr\'olikowski~\autocite{FK87}
compared the Petlyuk column with side strippers, side rectifiers and the direct and indirect sequences in closed form;
global MINLP studies~\autocite{TGAT,Nallasivam2016} find the Petlyuk value optimal on the instances they solve; we found
no proof (novelty checks in \texttt{references.md}).

\emph{Version~1 of this paper} (independently verified) proved the bound in Underwood's infinite-stage model for networks
of two-section columns with interval streams; the verifier found that the model's permissive form admits a
five-component point below the Petlyuk value, and listed what remained open: several feed stages, stream splitters,
recycles, reverse net flows and degenerate operating points. Version~1 is kept as
Sections~\ref{sec:setting}--\ref{sec:achieve}.

\emph{Version~2} replaces the Underwood model by the physical model itself: constant-molar-overflow networks of
equilibrium stages with any number of stages, of which ``infinitely many stages and pure products'' is a limit. In
this model the lower bound has an elementary proof that needs no information on sections, roots or topology.

\begin{theorem}[main results, version 2]\label{thm:main2}
Let $\alpha_1>\dots>\alpha_N>0$, $z_p>0$, $\sum_pz_p=1$, $q\in\mathbb R$, let $\theta_j\in(\alpha_{j+1},\alpha_j)$ be the roots of
$\sum_p\alpha_pz_p/(\alpha_p-\theta)=1-q$ and $H_j=\sum_{p>j}\alpha_pz_p/(\theta_j-\alpha_p)$.
\begin{enumerate}
\item (Lower bound, every arrangement.) At every steady state of every arrangement of the class $\mathcal A$
(Definition~\ref{def:A}: any finite network of equilibrium stages, single-phase splitters and mixers, total condensers
and total vaporizers, one feed split and conditioned at will, heat duties at any stage, recycles, products of either
phase), with any numbers of stages and any product purities, the total heat supplied $R$ satisfies
$R\ge\sum_P\max(0,\varphi_P(\theta_j))$ for every $j$, where $\varphi_P$ is the Underwood flux of the product stream $P$
(Theorem~\ref{thm:stage}); if every product has purity at least $1-\varepsilon$, then $R\ge\max_j(H_j-c_j\varepsilon)$ with
an explicit $c_j$ (Corollary~\ref{cor:purity}). Hence every arrangement needs, in the limit of infinitely many stages and
pure products, at least $\max_jH_j$.
\item (Attainment.) The generalized Petlyuk arrangement attains $\max_jH_j$ in Underwood's infinite-stage model
(Theorem~\ref{thm:achieve}, version~1). In the finite-stage model a single column attains $H_1$ for $N=2$
(Proposition~\ref{prop:binary}); for $N\ge3$ the finite-stage Petlyuk minimum approaches $\max_jH_j$ in our simulations
(Section~\ref{sec:attain}). Finite-stage attainment of the closed Petlyuk network for $N\ge3$ is not proved here.
\end{enumerate}
Consequently no arrangement of $\mathcal A$ needs less total vaporization than the minimum vapor $\max_jH_j$ of the
generalized Petlyuk arrangement computed by Halvorsen and Skogestad, the highest peak of the $V_{\min}$ diagram: the
conjecture holds in $\mathcal A$ in this sense, for every $N$. For $N=2$ it holds with the Petlyuk value also realized in
the finite-stage model; for $N\ge3$ the realization of $\max_jH_j$ as a finite-stage limit by the Petlyuk arrangement
(Halvorsen and Skogestad's formula read in that sense) is supported numerically but not proved here
(Section~\ref{sec:attain}).
\end{theorem}

The class $\mathcal A$ contains every arrangement Halvorsen and Skogestad have in mind and more: non-adiabatic sections
with external utilities, Murphree trays, any recycles. The proof (Section~\ref{sec:stage}) rests on one observation:
the Underwood flux $\varphi(\theta)=v-\sum_i\alpha_in_i/(\alpha_i-\theta)$ of every stream leaving an equilibrium stage has
the sign of $-\sum_i\alpha_ix_i/(\alpha_i-\theta)$, $x$ the stage liquid, whatever the phase. At $\theta=\theta_j$ the units
therefore split into a light side and a heavy side, and the energy balance of the heavy side, in Underwood's weights,
gives the bound. It was found independently by three agents of version~2 (Section~\ref{sec:credits2}).

\section{The class \texorpdfstring{$\mathcal A$}{A} and the physical model}\label{sec:classA}

The sources fix the physical assumptions: constant relative volatilities, constant molar flows, constant pressure,
equilibrium stages, one feed of arbitrary composition and thermal state, pure products, and the total vaporization
(heat input) as objective~\autocite{HS3,HS02,HalvorsenThesis}. Halvorsen's thesis (\S6.1.1) describes an arrangement as
sections of connected equilibrium stages, reboilers and condensers, directly coupled by two-way liquid and vapor
streams or coupled via condensers or reboilers, divided at stages with feed or side-draw streams; under constant molar
overflow the energy balance of any stage reduces to ``vapor out minus vapor in equals heat input'' (\emph{ibid.},
\S2.6.1). We take the molar heat of vaporization as the unit of heat.

\begin{definition}[the class $\mathcal A$]\label{def:A}
An \emph{arrangement} is a finite directed multigraph whose vertices are \emph{units} and whose arcs are
\emph{streams}, together with a feed source and product outlets. A stream $e$ has a phase (saturated liquid or
saturated vapor), component flows $n_{e,i}\ge0$ and vapor flow $v_e$ ($=\sum_in_{e,i}$ for a vapor, $0$ for a liquid).
The units are of four types:
\begin{itemize}
\item[(E)] \emph{equilibrium stage}: any inflows; an external heat duty $Q_u\in\mathbb R$; its contents leave as a liquid of
composition $x_u$ and a vapor of composition $y_u$ with $y_{u,i}=\alpha_ix_{u,i}/S(x_u)$, $S(x)=\sum_k\alpha_kx_k$, each
divided into any number of streams;
\item[(M)] \emph{single-phase mixer or splitter}: inflows of one phase, outflows of that phase with the mixed
composition, $Q_u=0$;
\item[(C)] \emph{total condenser}: vapor inflows, liquid outflows of the mixed composition, $Q_u=-\sum_{\mathrm{in}}v_e$;
\item[(R)] \emph{total vaporizer}: liquid inflows, vapor outflows of the mixed composition, $Q_u=+\sum_{\mathrm{in}}\sum_in_{e,i}$.
\end{itemize}
The process feed ($F=1$, composition $z$, liquid fraction $q$) is divided into portions of sizes $f_u\ge0$, $\sum f_u=1$,
entering units; every stream leaving the arrangement is a \emph{product stream}, assigned to one component. A
\emph{steady state} satisfies, at every unit, the component balances $\sum_{\mathrm{out}}n_{e,i}-\sum_{\mathrm{in}}n_{e,i}=f_uz_i$,
the energy balance $\sum_{\mathrm{out}}v_e-\sum_{\mathrm{in}}v_e=Q_u+(1-q)f_u$, and the composition rule of its type. Its
\emph{total vaporization} is $R=\sum_u\max(Q_u,0)$, the total heat supplied.
\end{definition}

Nothing is assumed on the topology (recycles are cycles of the graph), on the signs of the net component flows in a
section, on which components are present where, on the numbers of stages, or on the stability or uniqueness of the
steady state. An \emph{adiabatic section} is a chain of (E) units with $Q=0$, liquid flowing down and vapor flowing up
with constant flows; a column is a stack of sections joined at junction stages. Every device of the arrangements
considered in~\autocite{HS3,HalvorsenThesis,TGAT,SA10,Fidkowski2006,Agrawal2003} is a sub-network of units of
$\mathcal A$ (Agent~N's table, \texttt{total/network/network.md} \S1.2): conventional columns with total or partial
condensers and reboilers; thermal couplings; side draws of either phase; side strippers and rectifiers; Petlyuk,
generalized Petlyuk and Kaibel arrangements and dividing-wall columns (a liquid or vapor split is a stage outflow divided
into several streams, a merge is several inflows into one stage); several feed stages per column and intermediate
sections; stream splitters dividing a section product among several columns, remixing of any streams, vapor or liquid
bypasses, recycles; feed splitting, preheating, precooling and pre-flashing; intermediate reboilers and condensers,
heaters and coolers at any stage (non-adiabatic sections with external utilities, which the conjecture's wording does
not require us to admit); and trays with a Murphree vapor efficiency $E\in[0,1]$ common to all components, which are
exactly an (E) unit receiving the fraction $E$ of the vapor from below, a vapor bypass of the fraction $1-E$ and an (M)
mixer ($Lx_{\mathrm{in}}+EVy_{\mathrm{in}}=Lx+EVK(x)x$), hence also packed sections as limits.

\begin{definition}[infinitely many stages and pure products]\label{def:limit}
A steady state has \emph{purity} at least $1-\varepsilon$ if every product stream assigned to component $p$ has mole
fraction of $p$ at least $1-\varepsilon$. For a family $\mathcal S\subseteq\mathcal A$ (for example all stage numbers of one
configuration) let $\Sigma_\varepsilon(\mathcal S)$ be the set of steady states of members of $\mathcal S$ with purity at
least $1-\varepsilon$, and put
\[R^*(\mathcal S;\varepsilon)=\inf_{s\in\Sigma_\varepsilon(\mathcal S)}R(s),\qquad R^*(\mathcal S)=\lim_{\varepsilon\to0^+}R^*(\mathcal S;\varepsilon).\]
\end{definition}
The limit exists ($R^*(\mathcal S;\varepsilon)$ is nonincreasing in $\varepsilon$). Halvorsen and Skogestad's minimum vapor of a
configuration $X$ with infinitely many stages is $R^*(\mathcal S_X)$, $\mathcal S_X$ the arrangements with the topology of $X$:
every real arrangement has finitely many stages and impure products, and Underwood's values are the limits. In this
language the conjecture reads: $R^*(\mathcal S_{\mathrm{Petlyuk}})\le R^*(\mathcal S)$ for every $\mathcal S\subseteq\mathcal A$.

\emph{Excluded, and why it is outside the conjecture.} Heat exchange between process streams that is not counted in $R$
(``internal heat exchange''; \autocite{HS3}, \S6.3 and \S8; it lowers the external heat below the Petlyuk value,
\autocite{HalvorsenThesis}, Table~6.1); several pressures (excluded by constant pressure; they act through heat
exchange between columns or compression, \emph{ibid.}, \S6.6, and change $\alpha$); non-ideal mixtures and non-CMO
enthalpies; devices releasing streams of different compositions not in phase equilibrium (decanters, membranes,
component-dependent efficiencies); heat pumps; several process feeds of different compositions or thermal states.
Finite stages and impure products are not excluded: they are the objects of the model.

\section{The stage-level bound}\label{sec:stage}

For $\theta>0$, $\theta\notin\{\alpha_i\}$, the \emph{Underwood flux} of a stream $e$ and of the unit feed are
\[\varphi_e(\theta)=v_e-\sum_i\frac{\alpha_in_{e,i}}{\alpha_i-\theta},\qquad\varphi_F(\theta)=(1-q)-\sum_i\frac{\alpha_iz_i}{\alpha_i-\theta},\]
so that $\varphi_F(\theta_j)=0$ for every $j$. For a section with vapor $V$ and upward net flows $w$, the difference
$\varphi(\text{vapor up})-\varphi(\text{liquid down})$ at any cross-section is $V-\sum_i\alpha_iw_i/(\alpha_i-\theta)$, the
flux $\Lambda_s$ of version~1. Write $G_x(\theta)=\sum_i\alpha_ix_i/(\alpha_i-\theta)$.

\begin{lemma}[node balance]\label{lem:node}
At every unit $u$ of a steady state and every $\theta$, product streams counting as outflows,
\[\sum_{e\text{ out of }u}\varphi_e(\theta)-\sum_{e\text{ into }u}\varphi_e(\theta)=Q_u+f_u\varphi_F(\theta).\]
\end{lemma}
\begin{proof}
$\varphi$ is linear in $(v,n)$: $\varphi_e=v_e-\sum_ic_in_{e,i}$ with $c_i=\alpha_i/(\alpha_i-\theta)$. Subtract $\sum_ic_i$ times
the component balances from the energy balance of Definition~\ref{def:A}.
\end{proof}

\begin{lemma}[sign uniformity]\label{lem:uniform}
Let $\theta>0$, $\theta\notin\{\alpha_i\}$. A liquid stream of flow $l$ and composition $x$ has $\varphi=-l\,G_x(\theta)$; a vapor
stream of flow $v$ in equilibrium with a liquid of composition $x$ has $\varphi=-(\theta v/S(x))\,G_x(\theta)$. Hence all
streams leaving an (E) unit have fluxes of one sign, that of $-G_{x_u}(\theta)$, and so do all streams leaving an (M), (C) or
(R) unit (one phase, one composition).
\end{lemma}
\begin{proof}
For the vapor, with $\sum_iy_i=1$ and $y_i=\alpha_ix_i/S$: $\varphi=v\sum_iy_i\bigl(1-\frac{\alpha_i}{\alpha_i-\theta}\bigr)
=-\theta v\sum_i\frac{y_i}{\alpha_i-\theta}=-\frac{\theta v}{S}\sum_i\frac{\alpha_ix_i}{\alpha_i-\theta}$, and $\theta>0$, $S>0$.
\end{proof}

Fix $j$ and $\theta=\theta_j$. Call a unit \emph{heavy-ish} if all its outflows have $\varphi\ge0$ and \emph{light-ish}
otherwise; by Lemma~\ref{lem:uniform} the outflows of a light-ish unit all have $\varphi<0$ (or zero flow). Let
$\mathcal U_+$ and $\mathcal U_-$ be the two sets, $R(\mathcal W)=\sum_{u\in\mathcal W}\max(Q_u,0)$ and
$Q(\mathcal W)=\sum_{u\in\mathcal W}\max(-Q_u,0)$.

\begin{theorem}[stage-level cut]\label{thm:stage}
At every steady state of every arrangement of $\mathcal A$ and for every $j$,
\begin{equation}\label{eq:defect}
R-\sum_P\max\bigl(0,\varphi_P(\theta_j)\bigr)=R(\mathcal U_-)+Q(\mathcal U_+)+\sum_{e\colon\mathcal U_+\to\mathcal U_-}\varphi_e(\theta_j)
-\sum_{e\colon\mathcal U_-\to\mathcal U_+}\varphi_e(\theta_j),
\end{equation}
the sums running over the streams between the two sets, and every term on the right is nonnegative. In particular
$R\ge\sum_P\max(0,\varphi_P(\theta_j))$.
\end{theorem}
\begin{proof}
Sum Lemma~\ref{lem:node} over $u\in\mathcal U_+$ at $\theta_j$, where the feed terms vanish. Streams with both ends in
$\mathcal U_+$ cancel, so
$R(\mathcal U_+)-Q(\mathcal U_+)=\sum_{P\text{ from }\mathcal U_+}\varphi_P+\sum_{\mathcal U_+\to\mathcal U_-}\varphi_e-\sum_{\mathcal U_-\to\mathcal U_+}\varphi_e$.
Products leaving $\mathcal U_+$ have $\varphi_P\ge0$ and products leaving $\mathcal U_-$ have $\varphi_P\le0$, so the first sum is
$\sum_P\max(0,\varphi_P)$; streams leaving $\mathcal U_+$ have $\varphi\ge0$ and streams leaving $\mathcal U_-$ have $\varphi\le0$, so the
crossing terms are nonnegative. Add $R(\mathcal U_-)+Q(\mathcal U_+)$ to both sides.
\end{proof}

\begin{corollary}[pure products]\label{cor:pure}
If every product stream is pure, $R\ge H_j$ for every $j$; that is, $R\ge\max_jH_j$.
\end{corollary}
\begin{proof}
A pure liquid product of component $p$ and flow $P_p$ has $\varphi=\alpha_pP_p/(\theta_j-\alpha_p)$, a pure vapor product
$\varphi=\theta_jP_p/(\theta_j-\alpha_p)$; both are positive exactly for $p>j$, and the vapor value is the larger. The products
of component $p$ carry $z_p$ in total.
\end{proof}

\begin{corollary}[finite purity]\label{cor:purity}
If every product stream has purity at least $1-\varepsilon$, then for every $j$
\[R\ge H_j-c_j\varepsilon,\qquad c_j=\max\Bigl(\frac{\alpha_j}{\alpha_j-\theta_j},\frac{\alpha_{j+1}}{\theta_j-\alpha_{j+1}}\Bigr).\]
Consequently $R^*(\mathcal S)\ge\max_jH_j$ for every family $\mathcal S\subseteq\mathcal A$.
\end{corollary}
\begin{proof}
For a vapor stream, $\varphi^{\mathrm{vap}}=\varphi^{\mathrm{liq}}+(\text{flow})$ for the same component flows, so it suffices to
bound the liquid fluxes. Keep only the products assigned to components $p>j$:
$\sum_P\max(0,\varphi_P)\ge\sum_{i>j}\frac{\alpha_in^h_i}{\theta_j-\alpha_i}-\sum_{i\le j}\frac{\alpha_in^h_i}{\alpha_i-\theta_j}$,
$n^h_i$ the flow of $i$ in these products. By the overall balance $n^h_i=z_i-m_i$ for $i>j$, $m_i$ the flow of $i$ in
products assigned to components $\le j$. The flows $m_i$ ($i>j$) and $n^h_i$ ($i\le j$) are disjoint parts of the total
impurity, which is at most $\varepsilon$ times the total product flow $1$; the coefficients $\alpha_i/(\theta_j-\alpha_i)$ ($i>j$)
and $\alpha_i/(\alpha_i-\theta_j)$ ($i\le j$) are at most $c_j$.
\end{proof}

\noindent The proof uses only: constant molar overflow at every unit (Lemma~\ref{lem:node}); phase equilibrium with
constant relative volatilities, or a single outflow composition and phase (Lemma~\ref{lem:uniform}); $\theta_j>0$; one
feed composition and thermal state, so that every feed portion has zero flux at $\theta_j$; and the counting of every
positive duty in $R$. Agents F, S, N and P checked the identities and the bound on finite-stage steady states computed
by four independent solvers, and Agent N checked the identities in exact rational arithmetic
(Section~\ref{sec:computer2}).

\section{Consequences}\label{sec:conseq}

\emph{The Underwood pinch hyperplane.} For $0\le q\le1$ the equilibrium liquid $x^F$ of the feed satisfies
$G_{x^F}(\theta_j)=0$ for every $j$ (insert $z=(1-q)y^F+qx^F$ in the feed equation: $G_{x^F}(\theta_j)\,[(1-q)\theta_j/S(x^F)+q]=0$).
So a unit is light-ish or heavy-ish according to the side of the hyperplane $\{x:G_x(\theta_j)=0\}$, through the feed-flash
liquid, on which its liquid lies (for a (C) unit its liquid outflow; for an (R) unit the liquid in equilibrium with its
vapor outflow, since $\sum_iy_i/(\alpha_i-\theta)=G_x(\theta)/S(x)$ for $y$ in equilibrium with $x$). Theorem~\ref{thm:stage} says: \emph{the heat
supplied to the heavy side is at least $H_j$} (for pure products), and the excess of $R$ over the bound is the sum of three nonnegative
losses, heat supplied on the light side, heat removed on the heavy side, and flux across the hyperplane, as in pinch
analysis of heat-exchanger networks.

\emph{Sections.} In an adiabatic section, by Lemma~\ref{lem:uniform} applied to two consecutive stages,
$\Lambda_s(\theta)=L\,G_{x_{k+1}}(\theta)-(\theta V/S(x_k))\,G_{x_k}(\theta)$ at every cross-section (Underwood's identity, eq.~(2.49) of
\autocite{HalvorsenThesis} at the roots of $\Lambda_s$, here for every $\theta$). Hence along a section the side changes at
most once, from heavy-ish below to light-ish above if $\Lambda_s(\theta_j)>0$ and the reverse if $\Lambda_s(\theta_j)<0$, and a
section whose stages change side contributes $|\Lambda_s(\theta_j)|$ to the defect~\eqref{eq:defect}. At a pinch composition
$x^*\propto(w_i/(\alpha_i-\phi))_i$ of a section ($\phi$ a root of $\Lambda_s$), $G_{x^*}(\theta)=-(S(x^*)/V)\Lambda_s(\theta)/(\theta-\phi)$.
(\texttt{total/sections/sections.md}, Lemmas~6--9; \texttt{total/network/network.md}, Corollary~6.2.)

\emph{Ties.} By~\eqref{eq:defect}, a family of steady states with purity $\to1$ and $R\to H_j$ must have, in the limit, no heat
supplied to light-ish units, no heat removed from heavy-ish units, zero flux on every stream between the two sides, and
$\theta_j$ as an Underwood root of every section whose stages change side, with the change taking place in a zone where the
liquid is on the hyperplane: Underwood's pinch at an active root, now derived for every arrangement. In particular, at
an optimum with a strict highest peak $j^*$, no reboiler or heater may deliver a light product $p\le j^*$ and no condenser
or cooler a heavy product $p>j^*$ (Agent N, \texttt{network.md}, Corollary~5.3; version~1's Corollary~\ref{cor:tie} in
the whole class). With finitely many stages the inequality is strict whenever some section changes side.

\emph{What the bound covers.} Theorem~\ref{thm:stage} contains no hypothesis on reverse net flows (Halvorsen's regimes
in which a component flows against the direction of a directly coupled section~\autocite{HalvorsenThesis}, \S4.3.3),
degenerate distributions, several feed stages, intermediate sections, stream splitters, remixing, recycles, bypasses,
non-interval streams, products keeping every component of a column feed, or products of either phase: every gap of
version~1 listed in its verification is closed for the lower bound. It also bounds the total vaporization of
non-adiabatic arrangements with external utilities, such as reversible columns (whose total heat supply indeed exceeds
the adiabatic Petlyuk value in Halvorsen's examples~\autocite{HalvorsenThesis}, Table~6.1), and of Murphree trays.

\emph{Internal heat exchange and several pressures.} The theorem bounds the heat delivered to the heavy-ish units,
whatever its origin. With heat exchange between process streams counted in $R$ the bound still holds; the external
heat can only fall below $\max_jH_j$ by moving heat into heavy-ish units from light-ish ones, which is what
Halvorsen's heat-integrated examples do (side-stream condensation heating the feed, heat transfer across the dividing
wall, double-effect columns). This is why the conjecture must exclude internal heat exchange, and it answers the
question left in~\autocite{HS02} for the total vaporization.

\emph{The permissive counterexample of version~1 is not physical.} By Corollary~\ref{cor:purity}, any family of steady
states of configuration 161 (or of any arrangement) with purity $\to1$ has $\liminf R\ge\max_jH_j$, while the permissive
point of Section~\ref{sec:cex} has $R=H_3-2.6\cdot10^{-4}$. Concretely, its Underwood fluxes at $\theta_3$ contain a
light-to-heavy sign path $\mathrm C\to\mathrm{CDE}\to\mathrm{BCDE}\to\mathrm{BCD}\to\mathrm D$, impossible at a steady state because
along every arc of the sign graph a light-ish node can only be followed by a light-ish node
(\texttt{total/network/network.md}, \S6.3--6.4). The violation is global (a chain of four sections), which is why no
column-local argument could exclude it. For global optimization this suggests a repair of the permissive model that
keeps every physically admissible excess vapor: require, for every $j$, junction signs ($+$ at light products, $-$ at
heavy products) compatible with the signs of the section and link fluxes (one binary variable per junction and peak).
Every point of the repaired model satisfies $R\ge H_j$, and limits of physical steady states satisfy its constraints
when each junction of the Underwood description is one equilibrium stage (\emph{ibid.}, \S6.5).

\emph{Version~1's strict theorem.} The sign-graph argument of Theorem~\ref{thm:strict}, which needed the column
Lemmas~\ref{lem:sign} and~\ref{lem:interior} in the strict Underwood model, holds at every physical steady state with the
partition $U=\mathcal U_-$: the node sign is a potential that never decreases along an arc of the sign graph.

\section{Attainment}\label{sec:attain}

By Corollary~\ref{cor:purity} applied to the Petlyuk arrangement itself, $R^*(\mathcal S_{\mathrm{Petlyuk}})\ge\max_jH_j$; this is
the part of Halvorsen and Skogestad's formula for the Petlyuk minimum that they proved only for the extreme peaks, now
for every peak, every $N$ and every $q$ in the physical model. Since every arrangement satisfies the same bound, the
conjecture $R^*(\mathcal S_{\mathrm{Petlyuk}})\le R^*(\mathcal S)$ for all $\mathcal S$ \emph{follows from} the equality
$R^*(\mathcal S_{\mathrm{Petlyuk}})=\max_jH_j$ (\emph{attainment}: finite-stage Petlyuk steady states with purities tending to $1$
and $R$ tending to $\max_jH_j$), which is Halvorsen and Skogestad's formula for the Petlyuk minimum read in the
finite-stage-limit sense. (Attainment is sufficient, not logically necessary: without it the conjecture could still
hold, but the formula would fail.) In Underwood's infinite-stage model the generalized Petlyuk arrangement has an operating point with
$R=\max_jH_j$, every internal column at its preferred split (Theorem~\ref{thm:achieve} of version~1, independently
verified, for every $N$ and $q$; this is Halvorsen and Skogestad's construction~\autocite{HS3}). In the finite-stage
model:

\begin{proposition}[two components; Agent S]\label{prop:binary}
Let $N=2$, $\alpha_1=\alpha>1=\alpha_2$, light feed fraction $z\in(0,1)$, any $q$. For every $s>H_1$ and $\varepsilon>0$ there is $n_0$
such that a column with $n_T,n_B\ge n_0$ stages, a partial reboiler of duty $s$ and a total condenser has a steady
state with both products of purity at least $1-\varepsilon$. Hence $R^*(\mathcal S_{\mathrm{column}})=H_1$, and for $N=2$ the
conjecture holds in $\mathcal A$.
\end{proposition}
\begin{proof}
Write $e(x)=\alpha x/(1+(\alpha-1)x)$ (light vapor fraction in equilibrium with light liquid fraction $x$), $D=z$,
$B=1-z$, $V_B=s$, $V_T=s+1-q$, $L_T=V_T-D$, $L_B=s+B$; $L_T>0$ because $H_1+1-q=\alpha z/(\alpha-\theta_1)>z$. The map
$\beta_0(x)=(s/L_B)e(x)$ (bottom section, pure bottom product) has the fixed points $0$ and
$\hat x_B=(\alpha s/L_B-1)/(\alpha-1)\in(0,1)$ and $\beta_0(x)>x$ on $(0,\hat x_B)$; the map $\tau_0(x)=e^{-1}((L_Tx+D)/V_T)$ (top
section going down, pure top product) has the fixed points $1$ and $\hat x_T=(V_T/L_T-1)/(\alpha-1)$ and $\tau_0(x)<x$ on
$(\hat x_T,1)$. Now $\hat x_T<\hat x_B$ iff $f(s)=\alpha s(s+1-q-z)-(s+1-q)(s+1-z)>0$; after elimination of $q$ by the feed
equation, $f$ is a quadratic with positive leading coefficient $\alpha-1$ and roots $H_1$ and
$s_-=1/(\alpha-1)-z\theta_1/(\alpha-\theta_1)<H_1$ (\texttt{total/sections/scripts/check\_binary\_attainment.py}, sympy), so
$\hat x_T<\hat x_B$ for $s>H_1$. For $\eta\in[0,1]$ put $x_B=\eta$ and $x_D(\eta)=1-(1-z)\eta/z$ (light balance). Iterating
$x_{k+1}=(s\,e(x_k)+B\eta)/L_B$ upward from the reboiler liquid $\eta$ ($n_B+1$ steps) gives a feed-stage liquid
$x^{\mathrm{bot}}_f(\eta)$; iterating downward from the top vapor $x_D(\eta)$, $x=e^{-1}(y)$, $y_{\mathrm{next}}=(L_Tx+Dx_D)/V_T$
($n_T+1$ stages) gives $x^{\mathrm{top}}_f(\eta)$. Both are continuous, all iterates stay in $[0,1]$, and equality
$x^{\mathrm{bot}}_f=x^{\mathrm{top}}_f$ yields a steady state (the feed-stage balance follows from the two cross-section balances
and the overall balance). At $\eta=0$, $x^{\mathrm{bot}}_f=0<1=x^{\mathrm{top}}_f$. Fix $\eta_1\in(0,1]$ with
$\eta_1\max(1,(1-z)/z)\le\varepsilon$. By monotonicity, $x^{\mathrm{bot}}_f(\eta_1)\ge\beta_0^{n_B+1}(\eta_1)\to\hat x_B$ (or $\ge\hat x_B$) and
$x^{\mathrm{top}}_f(\eta_1)\le\max(\hat x_T,\tau_0^{n_T}(e^{-1}(x_D(\eta_1))))\to\hat x_T$, so for large $n_T,n_B$,
$x^{\mathrm{bot}}_f(\eta_1)>x^{\mathrm{top}}_f(\eta_1)$, and the intermediate value theorem gives $\eta^*\in(0,\eta_1)$ with a steady
state of purities $1-\eta^*$ and $1-(1-z)\eta^*/z$, both $\ge1-\varepsilon$, and $R=s$.
\end{proof}

For $N\ge3$ the difficulty is not existence: steady states exist for every prescribed set of flows.

\begin{proposition}[existence; Agent N]\label{prop:exist}
Let an arrangement of $\mathcal A$ and the flows of all its streams and feed portions be prescribed, with the total mass
balance at every unit, every unit reachable from a unit receiving feed and a product outlet reachable from every
unit. Then for every $\alpha_1>\dots>\alpha_N>0$ a steady state exists, and every steady state has all mole fractions
in $[c,1]$ for some $c>0$ depending only on the flows, $z$ and $\alpha_1/\alpha_N$. In particular $R$ is determined by the
flows, and $R^*(\mathcal S)\le\rho$ as soon as, for every $\eta>0$, some stage numbers and flows with $R\le\rho+\eta$ admit
a steady state with all purities at least $1-\eta$.
\end{proposition}
\begin{proof}
(\texttt{total/network/network.md}, \S10.) An (M) or (C) unit is an (E) unit without vapor outflow and an (M) vapor
unit or an (R) unit one without liquid outflow, so the unknowns are one composition per unit. For fixed
$K_{k,i}=\alpha_i/S(x_k)$ the balances of component $i$ form a linear system whose matrix is a Z-matrix with positive
diagonal, weakly chained diagonally dominant by columns (the excess of column $k$ is the flow that unit $k$ sends to
product outlets, and every column is chained to such a unit), hence a nonsingular M-matrix with nonnegative inverse,
positive between units joined by a path. Summing over $i$ shows that $\sum_ix_{k,i}=1$ holds automatically at every
solution of the unnormalized system. Since $K_{k,i}\in[\alpha_N/\alpha_1,\alpha_1/\alpha_N]$, the positive entries of the
inverses are bounded below, which gives the a priori bounds, uniformly along the homotopy $\alpha(t)=(\alpha_i^t)$; at
$t=0$ the system is linear with a unique nondegenerate solution, so the Brouwer degree is nonzero for every $t$.
\end{proof}

So attainment amounts to the \emph{purity} of steady states at fixed flows. Isolated columns are understood. For
$N=2$ this is Proposition~\ref{prop:binary}. For a directly coupled prefractionator with any number $m\ge3$ of
components and any feed quality, operated with its preferred net split deformed so as to keep its internal common
roots and with a vapor margin $\eta$, there are exact, strictly positive finite-stage profiles whose unwanted extreme
components at the two returns are at most $C\eta^{-1}\rho^n$, $0<\rho<1$; for $m=3$ the major-component ratios of the two
returns can moreover be prescribed (\texttt{total/attainment-astra/attainment.md}, Theorems~1--2, GPT-6-Astra; re-derived and
re-implemented by Agent S, \texttt{total/sections/sections.md} \S8.5). There the returns are boundary data: closing them through the finishing columns
of the Petlyuk network, with every impurity controlled, is the missing step. The obstruction to a general proof is the coupling
of pinch zones through the thermal couplings: in the infinite-stage limit, every section that carries excess vapor and
is pinched at both ends has a composition front whose position is fixed only by exponentially small terms, so the
pinch-matching limit system is not transversal there (Agent N's computation of the steady-state Jacobian:
smallest singular value decreasing geometrically with $n$; Agent S's analysis of doubly pinched sections).

\emph{Numerical evidence.} With the flows of Theorem~\ref{thm:achieve} plus a circulation $\delta$ from the reboiler to the
condenser spread over all sections, so that $R=\max_jH_j+\delta$ exactly, and $n$ stages per section, the product
impurities of the generalized Petlyuk arrangement fall as follows (Agent N, \texttt{total/network/scripts/ftc\_fixed\_flows.py}).
For $N=3$ ($\alpha=(4,2,1)$, $z_p=1/3$, $q=1$, $\delta=0.002$) the largest impurity is $1.4\cdot10^{-2}$, $1.1\cdot10^{-3}$,
$4.5\cdot10^{-7}$ and $1.3\cdot10^{-14}$ at $n=20,40,80,160$ (Newton residuals at most $2\cdot10^{-8}$); Agent P's independent
solver, with a circulation $\delta=0.01$ through both sections of every column, gives $8.1\cdot10^{-3}$, $1.1\cdot10^{-4}$ and
$9.3\cdot10^{-9}$ at $n=20,40,80$, a geometric decay of about $0.8$ per stage, and Agent F's solver, with the
prefractionator at its preferred split plus $0.02$ on both flows and $R=\max_jH_j+0.04$, gives $6.4\cdot10^{-3}$,
$6.7\cdot10^{-5}$ and $4.5\cdot10^{-9}$ at $n=20,40,80$ (\texttt{total/foundations/scripts/petlyuk\_attain.py}). For $N=4$ (one random instance,
$\delta=0.01$) N obtains $7.8\cdot10^{-2}$, $3.0\cdot10^{-3}$, $5.9\cdot10^{-5}$ and $2.1\cdot10^{-6}$ at $n=10,40,80,120$, with
Newton residuals of only $10^{-5}$ to $5\cdot10^{-7}$, so these values are indicative. Minimizing $R$ at fixed purity
(Agent P, \texttt{total/simulation/}), the optimum approaches the finite-purity bound of Theorem~\ref{thm:stage} as $n$
grows ($N=3$, purity $0.99$: $R-\text{bound}=2.9\cdot10^{-2}$, $3.1\cdot10^{-4}$, $3.0\cdot10^{-8}$ at $n=20,40,80$; Agent F's
separate solver and optimizer reproduce $R=1.372571$ at $n=20$), and the bound tends to $\max_jH_j$ as the purity tends
to $1$. With the same fixed-flow construction, Agent F's solver also gives purities tending to $1$ for $q=0$, $0.5$ and $1.3$
(largest impurity $2.2\cdot10^{-5}$, $1.7\cdot10^{-5}$ and $1.4\cdot10^{-4}$ at $n=40$, $R=\max_jH_j+0.04$).

\emph{Status and route for $N\ge3$.} For $N=3$, Agents S and N have reduced the closed network to a finite problem: at
fixed flows with margins, the sharp limit fixes all net flows, the prefractionator feed-stage composition (Astra's
vector) and the side-draw composition (pure B), but leaves the major-component ratios at the two coupling junctions
free (the two neutral modes); at finite $n$ these and the two trace net flows of the prefractionator (C upward, A
downward) are fixed by exponentially small balances, giving four unknowns, to be handled by a Poincar\'e--Miranda or
degree argument in exponentially rescaled variables with Astra's prefractionator theorem and binary shooting
arguments for the two finishing columns (\texttt{total/network/network.md} \S10, \texttt{total/sections/sections.md}
\S8.5--8.6). This argument is not complete; finite-stage attainment for $N\ge3$ therefore remains open, although every
computation supports it.
